% Proposed isolated B166, new checkpoint34 source. No earlier bytes changed.
\subsection*{The entire filling cohomology and its actual source subspace}
\statement{B.166}{Proved}{The entire filling-kernel cohomology is a left-coset space of the actual fixed subgroup, with the contact source as its exact closed subspace}
\label{stmt:B166}
Use B.165's actual groups and closures and put
\[
 \begin{gathered}
 \mathcal K=\mathcal K_a/\mathcal N,\qquad
 \mathcal G=\mathcal P/\mathcal N
       =\mathcal K\rtimes\widehat\sigma(\widehat J),\\
 \mathcal H=\overline{H_{\mathrm{src}}}
       =\mathcal I\rtimes\widehat\sigma(\widehat J),
 \qquad
 \mathcal I=\overline{I_{\mathrm{src}}}
       =\mathcal H\cap\mathcal K .
 \end{gathered}
\]
Here \(\widehat\sigma\) is the actual filled section of B.154;
\(\mathcal H,\mathcal I\) are the closed actual contact groups,
not abstract completions of their discrete kernels.
Identify \(\widehat S\) with its embedded image in \(\mathcal G\).
Then
\begin{equation}\label{eq:B166-fixed}
 \widehat{\Gamma_0}^{\,\delta}
     =\mathcal G^\delta
     =\mathcal K^\delta\rtimes\widehat\sigma(\widehat J),
\end{equation}
and there is a natural pointed homeomorphism
\begin{equation}\label{eq:B166-homeomorphism}
 \widehat S\backslash\mathcal G^\delta
   \xrightarrow{\ \sim\ }H^1(\Delta,\mathcal N),
 \qquad
 \widehat S(t\mathcal N)\longmapsto[t^{-1}\delta(t)].
\end{equation}
The left-coset convention in \eqref{eq:B166-homeomorphism} is
essential. Under B.165's identification of
\(H^1(\Delta,\mathcal K_a)\) with \(\widehat J/\widehat S\),
the natural map sends the class of \(g\) to
\begin{equation}\label{eq:B166-inversion}
                   \pi_J(g)^{-1}\widehat S,
\end{equation}
where \(\pi_J:\mathcal G\to\widehat J\) is the actual split projection.

The actual contact-source image in \eqref{eq:B166-homeomorphism}
is precisely the closed subspace
\begin{equation}\label{eq:B166-source-subspace}
 \widehat S\backslash\mathcal H
       \ \subseteq\ \widehat S\backslash\mathcal G^\delta.
\end{equation}
Consequently the following are equivalent:
\begin{equation}\label{eq:B166-equivalences}
 \begin{gathered}
 \widehat{\Gamma_0}^{\,\delta}=\mathcal H;\\
 H^1(\Delta,\mathcal N)
     \text{ equals the image of }\widehat S\backslash\mathcal H;\\
 \ker\!\left[
 H^1(\Delta,\mathcal N)\longrightarrow H^1(\Delta,\mathcal K_a)
 \right]=\partial(\mathcal I).
 \end{gathered}
\end{equation}
Here \(\partial\) is B.164's injective continuous connecting map.
These are equivalent formulations of the remaining local contact
gate; none is asserted to hold.

\begin{proof}
\emph{Actual quotients and fixed subgroups.}
The filled diagram of B.139 gives
\(\mathcal N\le\mathcal K_a\), and B.165 gives
\(\mathcal P=\mathcal K_a\rtimes\widehat\lambda(\widehat J)\).
Every generic \(\delta\)-defect of \(\widehat\lambda\) lies in
\(\mathcal R\le\mathcal N\). Its quotient section
\(\widehat\sigma(\widehat J)\) is therefore pointwise fixed.
The generic completed fixed subgroup is \(\widehat S\).
Its intersection with \(\mathcal N\) is trivial, because the split
source projection kills \(\mathcal N\). This justifies the embedded
\(\widehat S\) in the quotient.
Every fixed element of \(\widehat{\Gamma_0}\) maps into
\((\widehat{\Gamma_a})^\delta=\widehat J\).
The actual preimage computation of B.165 then places it in
\(\mathcal G\). Its split expression proves
\eqref{eq:B166-fixed}. This concerns the actual annular comparison,
not the full global icosahedral centralizer.

\emph{A surjective connecting map without a lift choice.}
B.165 proves that
\(H^1(\Delta,\mathcal K_a)\to H^1(\Delta,\mathcal P)\)
is constant at the basepoint. Thus every \(\mathcal N\)-cocycle
is a \(\mathcal P\)-coboundary. Define the compact closed lift set
\[
 \mathcal L=\{t\in\mathcal P:
                       t^{-1}\delta(t)\in\mathcal N\}.
\]
Its actual quotient map onto \(\mathcal G^\delta\) is surjective:
any lift of a fixed coset has its defect in \(\mathcal N\).
The continuous map
\[
 t\longmapsto[t^{-1}\delta(t)]\in H^1(\Delta,\mathcal N)
\]
is constant on the fibers, since replacing \(t\) by \(tn\) gives
\[
 (tn)^{-1}\delta(tn)
       =n^{-1}\bigl(t^{-1}\delta(t)\bigr)\delta(n)
       \qquad(n\in\mathcal N).
\]
It descends continuously through this compact quotient.
It is onto because every \(\mathcal N\)-cocycle is a
\(\mathcal P\)-coboundary. No continuous selection of lifts has
been assumed.

\emph{Exactly the left-coset fibers.}
Let \(g_i=t_i\mathcal N\in\mathcal G^\delta\) and
\(c_i=t_i^{-1}\delta(t_i)\). Equality of their cohomology classes
means that for some \(n\in\mathcal N\),
\[
                         c_2=n^{-1}c_1\delta(n).
\]
Thus \(t_1n\) and \(t_2\) have the same crossed defect, and
\[
 s=t_1nt_2^{-1}\in\mathcal P^\delta=\widehat S.
\]
Passing to the quotient gives \(s=g_1g_2^{-1}\).
Conversely, if \(g_1=sg_2\) for \(s\in\widehat S\), choose
the same actual fixed lift \(s\in\mathcal P\) and write
\(t_1=st_2n^{-1}\) with \(n\in\mathcal N\).
The same calculation yields the displayed cohomology equivalence.
These are precisely the left \(\widehat S\)-cosets.
The induced map in \eqref{eq:B166-homeomorphism} is a continuous
bijection. Its domain is compact; its target is Hausdorff because
the cocycle set and its compact-group orbit relation are closed.
It is therefore a homeomorphism.

\emph{The inverse coset on the target.}
For \(t\in\mathcal P\), put \(j=\widehat q_a(t)\) and \(l=j^{-1}\).
As in B.163, write
\[
 t=\delta\widehat\lambda(l)^{-1}k,\qquad k\in\mathcal K_a.
\]
Then
\[
 t^{-1}\delta(t)=k^{-1}X_l\delta(k),
 \qquad X_l=\delta\widehat\lambda(l)\widehat\lambda(l)^{-1}.
\]
Its image in \(H^1(\Delta,\mathcal K_a)\) is labelled
\(l\widehat S=j^{-1}\widehat S\), proving
\eqref{eq:B166-inversion}. This is independent of changing the
left \(\widehat S\)-representative.
For an explicit section check, take
\(t=\delta\widehat\lambda(l)^{-1}\). Its defect is literally
\(X_l\), and its filled image is \(\widehat\sigma(l)^{-1}\).
Equivalently,
\[
 c_{\widehat\lambda(l^{-1})}=X_l^{-1},
 \qquad [X_l^{-1}]=[X_l],
\]
where the second identity uses conjugator \(X_l\) and
\(\delta(X_l)=X_l^{-1}\):
\[
                  X_l^{-1}X_l\delta(X_l)=X_l^{-1}.
\]
This checks the inverse convention using the literal generic section.

\emph{The entire contact subspace and its basepoint fiber.}
B.150/B.154/B.164 give the actual closed subgroup
\(\mathcal H=\mathcal I\rtimes\widehat\sigma(\widehat J)\)
of \(\mathcal G^\delta\), containing \(\widehat S\).
Its left-coset space embeds as the closed subspace
\eqref{eq:B166-source-subspace}; compactness supplies closedness.
In the basepoint fiber of the natural map to
\(\widehat J/\widehat S\), this subspace is exactly
\(\partial(\mathcal I)\).
Indeed, write \(h=i\widehat\sigma(j)\in\mathcal H\).
The target is the basepoint exactly when \(j\in\widehat S\).
Left multiplication by \(\widehat\sigma(j)^{-1}\) then gives
\(\widehat\sigma(j)^{-1}i\widehat\sigma(j)\in\mathcal I\),
by source invariance. Conversely every \(i\in\mathcal I\)
lies in this fiber. The restriction of the connecting map to
\(\mathcal K^\delta\) is the injective \(\partial\) of B.164,
because \(\mathcal K_a^\delta=1\).

Equality of the two left-coset spaces implies equality of the groups:
if the coset of \(g\) is represented by \(h\in\mathcal H\), then
\(g=sh\in\mathcal H\), since \(s\in\widehat S\le\mathcal H\).
The common split \(\widehat J\)-component and the exact basepoint
fiber description give all equivalences in
\eqref{eq:B166-equivalences}.
The known \(\mathcal R\)-family supplies the entire target coset
space and its pointed section. It does not compute the remaining
basepoint fiber. In particular the known \(D_s\) contact supplies
a nonbase class in that fiber, so whole-map injectivity is false.
Nothing here proves B.164's cofinal finite-neighborhood containment.
\end{proof}

\begingroup\footnotesize\raggedright
The complete prior proof, both inverse-convention checks, and exact
source bindings are retained in
\path{work_CE/entire_filling_cohomology_contact_subspace33.txt}.
Its complete independent review is
\path{work_DE/review_entire_filling_cohomology_contact_subspace33.md}.
Every actual node and signed source translate remains in
\(\mathcal N\). The remaining local contact equality and the
separate global Main B transfer are open.
\par\endgroup
